\clearpage
\phantomsection
\chapter*{题目索引}
\addcontentsline{toc}{chapter}{题目索引}
\par\medskip\noindent\textbf{第1章}\par
\BookIndexEntry{ds-ex1-1}{例 1-1　值以外还要存什么 〖1.1-005〗}
\BookIndexEntry{ds-ex1-2}{例 1-2　循环由累计和决定何时停止 【2017】}
\BookIndexEntry{ds-ex1-3}{例 1-3　两层循环为什么仍是线性时间 【2022】}
\BookIndexEntry{ds-ch1-check}{自测与答案}
\par\medskip\noindent\textbf{第2章}\par
\BookIndexEntry{ds-ex2-1}{例 2-1　没有给存储方式时能否确定移动数 〖2.3-012〗}
\BookIndexEntry{ds-ex2-2}{例 2-2　如何一次扫描找到倒数第 \(k\) 个结点 【2009】}
\BookIndexEntry{ds-ex2-3}{例 2-3　把结点插到双链表结点之前 〖2.3-015〗}
\BookIndexEntry{ds-ex2-4}{例 2-4　数组循环左移怎样避免反复搬移 【2010】}
\BookIndexEntry{ds-ex2-5}{例 2-5　相同后缀是相同字符还是相同结点 【2012】}
\BookIndexEntry{ds-ex2-6}{例 2-6　主元素为什么还要再计数 【2013】}
\BookIndexEntry{ds-ch2-check}{自测与答案}
\par\medskip\noindent\textbf{第3章}\par
\BookIndexEntry{ds-ex3-1}{例 3-1　队列不会改变栈的输出次序 【2009】}
\BookIndexEntry{ds-ex3-2}{例 3-2　换了 front 含义，长度公式是否改变 〖3.2-006〗}
\BookIndexEntry{ds-ex3-3}{例 3-3　三对角矩阵的起止行不能一律算3个 【2016】}
\BookIndexEntry{ds-ex3-4}{例 3-4　先用对称性换到实际保存的半边 【2020】}
\BookIndexEntry{ds-ex3-5}{例 3-5　弹出次序怎样影响减法 【2018】}
\BookIndexEntry{ds-ch3-check}{自测与答案}
\par\medskip\noindent\textbf{第4章}\par
\BookIndexEntry{ds-ex4-1}{例 4-1　第一次失配后哪个指针改变 【2015】}
\BookIndexEntry{ds-ex4-2}{例 4-2　修正后最长能移动几位 【2024】}
\BookIndexEntry{ds-ch4-check}{自测与答案}
\par\medskip\noindent\textbf{第5章}\par
\BookIndexEntry{ds-ex5-1}{例 5-1　怎样从非叶结点数求叶结点数 【2010】}
\BookIndexEntry{ds-ex5-2}{例 5-2　任意形态的顺序存储需要预留多少 【2020】}
\BookIndexEntry{ds-ex5-3}{例 5-3　中序前驱和后继在哪一侧 【2024】}
\BookIndexEntry{ds-ex5-4}{例 5-4　二叉树隔两条链接，在森林中是什么关系 【2009】}
\BookIndexEntry{ds-ex5-5}{例 5-5　加权平均码长怎样从合并过程得到 【2023】}
\BookIndexEntry{ds-ex5-6}{例 5-6　已占用的码字怎样限制剩余容量 〖5.5-004〗}
\BookIndexEntry{ds-ch5-check}{自测与答案}
\par\medskip\noindent\textbf{第6章}\par
\BookIndexEntry{ds-ex6-1}{例 6-1　``有可能连通''与``保证连通''不同 【2010】}
\BookIndexEntry{ds-ex6-2}{例 6-2　什么时候允许回退去另一个分支 〖6.3-011〗}
\BookIndexEntry{ds-ex6-3}{例 6-3　同权边怎样影响第二次选边 【2015】}
\BookIndexEntry{ds-ex6-4}{例 6-4　求出第二条最短路径后哪些值还可变 【2021】}
\BookIndexEntry{ds-ex6-5}{例 6-5　怎样判定唯一拓扑序列 【2024】}
\BookIndexEntry{ds-ex6-6}{例 6-6　多条关键路径时怎样缩短工期 【2013】}
\BookIndexEntry{ds-ch6-check}{自测与答案}
\par\medskip\noindent\textbf{第7章}\par
\BookIndexEntry{ds-ex7-1}{例 7-1　两级顺序查找怎样选择块长 【2025】}
\BookIndexEntry{ds-ex7-2}{例 7-2　失败查找的最大比较次数 【2010】}
\BookIndexEntry{ds-ex7-3}{例 7-3　连续插入时在哪里旋转 〖7.3-003〗}
\BookIndexEntry{ds-ex7-4}{例 7-4　中序为降序时搜索方向怎样变 【2015】}
\BookIndexEntry{ds-ex7-5}{例 7-5　关键字总数怎样限制高度 〖7.4-009〗}
\BookIndexEntry{ds-ex7-6}{例 7-6　小规模B树的形态怎样完整计数 【2025】}
\BookIndexEntry{ds-ex7-7}{例 7-7　删除标记为什么继续计入失败路径 【2023】}
\BookIndexEntry{ds-ex7-8}{例 7-8　二次探查必须回到原始地址计算 【2024】}
\BookIndexEntry{ds-ch7-check}{自测与答案}
\par\medskip\noindent\textbf{第8章}\par
\BookIndexEntry{ds-ex8-1}{例 8-1　两趟希尔排序怎样按下标分组 〖8.2-012〗}
\BookIndexEntry{ds-ex8-2}{例 8-2　第二趟结束后应有几个枢轴就位 【2019】}
\BookIndexEntry{ds-ex8-3}{例 8-3　删除堆顶时怎样数关键字比较 【2015】}
\BookIndexEntry{ds-ex8-4}{例 8-4　比较次数不等于移动次数 【2024】}
\BookIndexEntry{ds-ex8-5}{例 8-5　第一趟后为什么尚未整体递增 〖8.5-012〗}
\BookIndexEntry{ds-ex8-6}{例 8-6　冻结记录在什么时候重新参加 【2023】}
\BookIndexEntry{ds-ex8-7}{例 8-7　不等长有序表怎样安排合并 【2012】}
\BookIndexEntry{ds-ex8-8}{例 8-8　四路最佳归并为什么补零权段 〖8.7-004〗}
\BookIndexEntry{ds-ch8-check}{自测与答案}
\clearpage
\phantomsection
\chapter*{方法索引}
\addcontentsline{toc}{chapter}{方法索引}
\par\medskip\noindent\textbf{第1章}\par
\BookIndexEntry{ds-01-01}{逻辑关系怎样落实到存储}
\BookIndexEntry{ds-01-02}{算法怎样保证结束并得到正确结果}
\BookIndexEntry{ds-01-03}{怎样由执行次数得到时间复杂度}
\BookIndexEntry{ds-c01-h41}{先找循环变量怎样变化}
\BookIndexEntry{ds-01-04}{辅助空间怎样包含递归栈}
\par\medskip\noindent\textbf{第2章}\par
\BookIndexEntry{ds-02-01}{位序、下标和基本操作}
\BookIndexEntry{ds-02-02}{顺序表的寻址、插入和删除}
\BookIndexEntry{ds-02-03}{单链表：先保留哪条链接}
\BookIndexEntry{ds-02-04}{双链表、循环链表与静态链表}
\BookIndexEntry{ds-02-05}{双指针、归并与原地处理}
\par\medskip\noindent\textbf{第3章}\par
\BookIndexEntry{ds-03-01}{栈的后进先出与合法序列}
\BookIndexEntry{ds-03-02}{队列与双端队列}
\BookIndexEntry{ds-03-03}{顺序栈、共享栈与循环队列}
\BookIndexEntry{ds-c03-h47}{循环队列先立约定再计算}
\BookIndexEntry{ds-03-04}{链栈与链队列}
\BookIndexEntry{ds-03-05}{多维数组的地址怎样计算}
\BookIndexEntry{ds-03-06}{对称、三角、三对角与稀疏矩阵}
\BookIndexEntry{ds-03-07}{表达式、递归与层次处理}
\par\medskip\noindent\textbf{第4章}\par
\BookIndexEntry{ds-04-01}{串、子串与下标}
\BookIndexEntry{ds-04-02}{朴素匹配怎样回退}
\BookIndexEntry{ds-04-03}{KMP 的 next 为什么只依赖模式串}
\par\medskip\noindent\textbf{第5章}\par
\BookIndexEntry{ds-05-01}{树的结点、度、层次与路径}
\BookIndexEntry{ds-05-02}{二叉树、满二叉树与完全二叉树}
\BookIndexEntry{ds-05-03}{顺序编号与链式表示}
\BookIndexEntry{ds-05-04}{遍历、还原与树形递归}
\BookIndexEntry{ds-05-05}{线索二叉树：区分孩子链接与遍历线索}
\BookIndexEntry{ds-05-06}{普通树的三种存储}
\BookIndexEntry{ds-05-07}{树、森林怎样转换为二叉树}
\BookIndexEntry{ds-05-08}{树与森林的遍历怎样对应}
\BookIndexEntry{ds-05-09}{哈夫曼树、WPL 与前缀编码}
\BookIndexEntry{ds-05-10}{并查集：用根结点代表一个集合}
\BookIndexEntry{ds-05-11}{堆的结构、调整与优先队列}
\par\medskip\noindent\textbf{第6章}\par
\BookIndexEntry{ds-06-01}{顶点、边、连通性与路径}
\BookIndexEntry{ds-06-02}{邻接矩阵、邻接表与基本操作}
\BookIndexEntry{ds-06-03}{十字链表与邻接多重表}
\BookIndexEntry{ds-06-04}{深度优先搜索与搜索森林}
\BookIndexEntry{ds-06-05}{BFS 为什么得到无权最短路}
\BookIndexEntry{ds-06-06}{Prim 与 Kruskal 在挑什么边}
\BookIndexEntry{ds-06-07}{Dijkstra 与 Floyd 的条件和中间状态}
\BookIndexEntry{ds-06-08}{拓扑排序：由入度判断哪些任务可以先做}
\BookIndexEntry{ds-06-09}{关键路径：事件与活动的时间不同}
\par\medskip\noindent\textbf{第7章}\par
\BookIndexEntry{ds-07-01}{关键字、成功失败与平均查找长度}
\BookIndexEntry{ds-07-02}{顺序查找、分块查找与提前结束}
\BookIndexEntry{ds-07-03}{折半查找与判定树}
\BookIndexEntry{ds-07-04}{二叉搜索树的区间不变式}
\BookIndexEntry{ds-07-05}{AVL：先定位最小失衡子树}
\BookIndexEntry{ds-07-06}{红黑树：颜色约束怎样限制高度}
\BookIndexEntry{ds-07-07}{B树：一个结点保存多个分隔关键字}
\BookIndexEntry{ds-07-09}{散列表：冲突、删除与ASL}
\BookIndexEntry{ds-07-10}{怎样按操作约束选择方法}
\par\medskip\noindent\textbf{第8章}\par
\BookIndexEntry{ds-08-01}{稳定性、内外排序与比较口径}
\BookIndexEntry{ds-08-02}{直接插入与折半插入}
\BookIndexEntry{ds-08-03}{起泡：一次相邻交换消除一个逆序}
\BookIndexEntry{ds-08-04}{简单选择：比较与移动的代价不同}
\BookIndexEntry{ds-08-05}{希尔：一趟只保证各增量子序列有序}
\BookIndexEntry{ds-08-06}{快速排序：先划分，再递归}
\BookIndexEntry{ds-08-07}{堆排序：堆范围与有序区分开}
\BookIndexEntry{ds-08-08}{二路归并：有序子段怎样变成长有序段}
\BookIndexEntry{ds-08-09}{基数排序：低位排好后怎样保留次序}
\BookIndexEntry{ds-08-10}{外部排序：归并段、缓冲与I/O}
\BookIndexEntry{ds-c08-h214}{败者树怎样减少每次输出后的比较}
\BookIndexEntry{ds-08-11}{怎样比较和选择排序算法}
